\section{数据引擎：从失败到成功的数据闭环}

谢晨说，最有效的数据是“先失败再成功”的数据。这句话很重要。成功轨迹告诉模型该怎么做，失败再修正的轨迹告诉模型为什么错、如何恢复、怎样从错误中学习。自动驾驶和机器人都需要这样的闭环：部署、采集、挖掘、训练、评测、再部署。

\begin{figure}[H]
\centering
\includegraphics[width=0.94\textwidth]{data-engine-loop.png}
\caption{数据闭环 / Data Engine：部署、采集、挖掘、训练、评测、仿真。}
\end{figure}

\begin{knowledgebox}{读图：Data Engine 的关键是持续性}
图中每一轮部署都会带回真实反馈，失败案例被挖掘成数据，训练后再通过仿真评测验证，最后重新部署。它支持的结论是：数据不是一次性购买，而是系统持续运行的结果。
\end{knowledgebox}

\subsection{特斯拉数据引擎的启示}

特斯拉的数据引擎来自大量端侧车辆。车在真实世界运行，产生长尾场景和驾驶反馈，再回到云端训练，更新端侧能力。这个逻辑在机器人上不完全成立，因为机器人端侧规模不足。于是，机器人产业必须寻找新的数据引擎：仿真、合成、遥操作、失败案例挖掘、benchmark 和真实小规模部署结合。

\begin{importantbox}{数据引擎的本质}
数据引擎不是“有很多数据”，而是有一个能持续发现错误、生成补充数据、验证改进并重新部署的循环。
\end{importantbox}

\subsection{本章小结}

高质量数据来自闭环。机器人和 LLM 都需要反馈，但机器人的反馈更昂贵、更依赖环境和评测。

